% TODO make hw stylesheet
\documentclass[10pt]{scrartcl}
\usepackage[a4paper, total={6in, 8in}]{geometry}
\usepackage[usenames,dvipsnames,svgnames]{xcolor}
\usepackage[shortlabels]{enumitem}
\usepackage[framemethod=TikZ]{mdframed}
\usepackage{amsmath,amssymb,amsthm}
\usepackage[colorlinks]{hyperref}
\usepackage{multicol}
\usepackage{tabularx}

\hypersetup{
	colorlinks=true,
	linkcolor=blue,
	filecolor=magenta,      
	urlcolor=magenta,
	pdftitle={MAT 322 HW}
}

\usepackage{microtype}
\usepackage{mathtools}
\usepackage[headsepline]{scrlayer-scrpage}
\usepackage{thmtools}
\usepackage[textsize=footnotesize]{todonotes}

\mdfdefinestyle{mdbluebox}{roundcorner=10pt,innerbottommargin=9pt,
	linecolor=blue,backgroundcolor=TealBlue!5,}
\declaretheoremstyle[headfont=\sffamily\bfseries\color{MidnightBlue},
mdframed={style=mdbluebox},]{thmbluebox}

\mdfdefinestyle{mdbloobox}{frametitlefont=\bfseries,innerbottommargin=8pt,
	nobreak=true,backgroundcolor=TealBlue!5,linecolor=blue,}
\declaretheoremstyle[headfont=\bfseries\color{MidnightBlue},
mdframed={style=mdbloobox},headpunct={\\[3pt]},postheadspace=0pt,]{thmbloobox}

\mdfdefinestyle{mdredbox}{frametitlefont=\bfseries,innerbottommargin=8pt,
	nobreak=true,backgroundcolor=Salmon!5,linecolor=RawSienna,}
\declaretheoremstyle[headfont=\bfseries\color{RawSienna},
mdframed={style=mdredbox},headpunct={\\[3pt]},postheadspace=0pt,]{thmredbox}

\mdfdefinestyle{mdgreenbox}{linecolor=ForestGreen,backgroundcolor=ForestGreen!5,
	linewidth=2pt,rightline=false,leftline=true,topline=false,bottomline=false,}
\declaretheoremstyle[headfont=\bfseries\sffamily\color{ForestGreen!70!black},
mdframed={style=mdgreenbox},headpunct={ --- },]{thmgreenbox}

\mdfdefinestyle{mdorangebox}{linecolor=RedOrange,backgroundcolor=RedOrange!5,
	linewidth=2pt,rightline=false,leftline=true,topline=false,bottomline=false,}
\declaretheoremstyle[headfont=\bfseries\sffamily\color{RedOrange!70!black},
mdframed={style=mdorangebox},headpunct={ --- },]{thmorangebox}

\mdfdefinestyle{mdblackbox}{linecolor=black,backgroundcolor=RedViolet!5!gray!5,
	linewidth=3pt,rightline=false,leftline=true,topline=false,bottomline=false,}
\declaretheoremstyle[mdframed={style=mdblackbox}]{thmblackbox}

\declaretheoremstyle[spaceabove=6pt,spacebelow=6pt]{basehead}

\declaretheorem[style=basehead,name=Problem]{problem}
\declaretheorem[style=thmredbox,name=Problem,numbered=no]{problem*}
\declaretheorem[style=thmbloobox,name=Setup]{setup} % for config geo
\declaretheorem[style=thmredbox,name=Example,sibling=problem]{example}
\declaretheorem[style=thmredbox,name=$\bigstar$ Problem,sibling=problem]{niceprob}
\declaretheorem[style=thmbluebox,name=Theorem,sibling=problem]{theorem}
\declaretheorem[style=thmbluebox,name=Corollary,sibling=theorem]{corollary}
\declaretheorem[style=thmbluebox,name=Proposition,sibling=theorem]{prop}
\declaretheorem[style=thmbluebox,name=Lemma,numberwithin=problem]{lemma}
\declaretheorem[style=thmbluebox,name=Theorem,numbered=no]{theorem*}
\declaretheorem[style=thmbluebox,name=Lemma,numbered=no]{lemma*}
\declaretheorem[style=thmgreenbox,name=Claim,numberwithin=problem]{claim}
\declaretheorem[style=thmgreenbox,name=Claim,numbered=no]{claim*}
\declaretheorem[style=thmblackbox,name=Remark,numberwithin=problem]{remark}
\declaretheorem[style=thmblackbox,name=Remark,numbered=no]{remark*}
\declaretheorem[style=thmgreenbox,name=Definition,sibling=theorem]{definition}
\declaretheorem[style=thmgreenbox,name=Definition,numbered=no]{definition*}
\declaretheorem[style=thmorangebox,name=Warning,sibling=theorem]{warning}
\declaretheorem[style=thmorangebox,name=Warning,numbered=no]{warning*}
\declaretheorem[style=thmorangebox,name=Note,numbered=no]{note}
\declaretheorem[style=thmblackbox,name=Exercise,sibling=problem]{exercise}
\declaretheorem[style=thmblackbox,name=Exercise,numbered=no]{exercise*}
\declaretheorem[style=thmblackbox,name=Question,sibling=problem]{question}
\declaretheorem[style=thmblackbox,name=Question,numbered=no]{question*}

\addtolength{\textheight}{3.14cm}
\providecommand{\ol}{\overline}
\providecommand{\ul}{\underline}
\providecommand{\eps}{\varepsilon}
\providecommand{\half}{\frac{1}{2}}
\providecommand{\dang}{\measuredangle} %% Directed angle
\providecommand{\CC}{\mathbb C}
\providecommand{\FF}{\mathbb F}
\providecommand{\NN}{\mathbb N}
\providecommand{\QQ}{\mathbb Q}
\providecommand{\RR}{\mathbb R}
\providecommand{\ZZ}{\mathbb Z}
\providecommand{\dg}{^\circ}
\providecommand{\ii}{\item}
\providecommand{\setdiff}{\smallsetminus}
\providecommand{\alert}[1]{{\sffamily\textbf{\textcolor{blue}{#1}}}}
\providecommand{\inv}{^{-1}}
\providecommand{\mb}[1]{\mathbf{#1}}
\providecommand{\dif}{\;d}

\reversemarginpar

\newenvironment{soln}{\begin{proof}[Solution]}{\end{proof}}
\newenvironment{walkthrough}{\noindent \textbf{Walkthrough.}}{\medskip}
\newlist{walk}{enumerate}{3}
\setlist[walk]{label=\bfseries (\alph*)}

\providecommand{\pow}{\operatorname{Pow}}
\providecommand{\vocab}[1]{{\textbf{\textcolor{ForestGreen}{#1}}}}
\providecommand{\lcm}{\operatorname{lcm}}
\providecommand{\abs}[1]{\left|#1\right|}

\usepackage{asymptote}
\begin{asydef}
	defaultpen(fontsize(10pt));
	unitsize(1cm);
	size(8cm); // set a reasonable default
	usepackage("amsmath");
	usepackage("amssymb");
	settings.tex="pdflatex";
	settings.outformat="pdf";
	// Replacement for olympiad+cse5 which is not standard
	import geometry;
	// recalibrate fill and filldraw for conics
	void filldraw(picture pic = currentpicture, conic g, pen fillpen=defaultpen, pen drawpen=defaultpen)
	{ filldraw(pic, (path) g, fillpen, drawpen); }
	void fill(picture pic = currentpicture, conic g, pen p=defaultpen)
	{ filldraw(pic, (path) g, p); }
	// some geometry
	pair foot(pair P, pair A, pair B) { return foot(triangle(A,B,P).VC); }
	pair orthocenter(pair A, pair B, pair C) { return orthocentercenter(A,B,C); }
	pair centroid(pair A, pair B, pair C) { return (A+B+C)/3; }
	// cse5 abbreviations
	path CP(pair P, pair A) { return circle(P, abs(A-P)); }
	path CR(pair P, real r) { return circle(P, r); }
	pair IP(path p, path q) { return intersectionpoints(p,q)[0]; }
	pair OP(path p, path q) { return intersectionpoints(p,q)[1]; }
	path Line(pair A, pair B, real a=0.6, real b=a) { return (a*(A-B)+A)--(b*(B-A)+B); }
	// cse5 more useful functions
	picture CC() {
		picture p=rotate(0)*currentpicture;
		currentpicture.erase();
		return p;
	}
	pair MP(Label s, pair A, pair B = plain.S, pen p = defaultpen) {
		Label L = s;
		L.s = "$"+s.s+"$";
		label(L, A, B, p);
		return A;
	}
	pair Drawing(Label s = "", pair A, pair B = plain.S, pen p = defaultpen) {
		dot(MP(s, A, B, p), p);
		return A;
	}
	path Drawing(path g, pen p = defaultpen, arrowbar ar = None) {
		draw(g, p, ar);
		return g;
	}
\end{asydef}

\usepackage{mathrsfs}

\ihead{\footnotesize MAT 322 HW06}
\ohead{\footnotesize Last compiled \today}

\title{\vspace{-2em}MAT 322 HW06}
\author{\large Aditya Pahuja}
\date{\large Due 04/01}
\begin{document}
\maketitle

\begin{problem}
	Let $S$ be the tetrahedron in $\RR^3$ having vertices $(0,0,0)$,
	$(1,2,3)$, $(0,1,2)$, and $(-1,1,1)$.
	Evaluate $\int_Sf$ where $f(x,y,z) = x+2y-z$.
\end{problem}

\begin{soln}
	Let $L$ be the linear transformation sending
	the tetrahedron $T$ formed by the vectors $(1,0,0)$, $(0,1,0)$, $(0,0,1)$
	to $S$, specifically
	\[L = \begin{pmatrix}
		1&0&-1\\2&1&1\\3&2&1
	\end{pmatrix}\]
	Then
	\[\int_{L(T)}f = \int_T (f\circ L)\cdot |\det(\mathrm dL)|.\]
	We can compute
	\[(f\circ L)(x,y,z) = f(x-z,2x+y+z,3x+2y+z)
	= (x-z) + 2(2x + y + z) - (3x + 2y + z)
	= 2x,\]
	and, since $L$ is linear, its derivative is equal to $L$,
	so $\det(\mathrm dL) = \det L = -2$. Thus we wish to compute
	\[4\int_T x.\]
	We can rewrite this as an iterated integral and compute
	\[4\int_{x=0}^1 \int_{y=0}^{1-x} \int_{z=0}^{1-x-y} x
	= \boxed{\frac16}.\qedhere\]
\end{soln}

\begin{problem}
	Let $U$ be the open set in $\RR^2$ consisting of all $x$
	with $\|x\|<1$. Let $f(x,y) = 1/(x^2+y^2)$ for $(x,y\ne(0,0))$.
	Determine whether $f$ is integrable over $U\setdiff\{0\}$
	and over $\RR^2\setdiff\ol U$. If so, evaluate the integral.
\end{problem}

\begin{soln}
	For $n>2$, let $K_n$ be the annulus formed by
	deleting the open disk centered at $(0,0)$
	with radius $1/n$ from the closed disk centered at $(0,0)$
	with radius $1-1/n$.
	Then, clearly $\bigcup_n K_n = U\setdiff \{0\}$.
	We can then integrate $\frac1{x^2+y^2}$ over $K_n$ via polar coordinates:
	\[\int_{K_n} f
	= \int_{[0,2\pi]\times[1/n,1-1/n]} \frac{1}{r^2}\cdot r
	= 2\pi(\ln(1-1/n) - \ln(1/n)) = 2\pi\cdot\ln(n-1).\]
	This is unbounded as $n$ grows large, so $f$
	is not integrable over $U$.
	
	Similarly, for $n > 2$, let $L_n$ be the annulus formed by
	deleting the open disk centered at $(0,0)$ with radius $1+1/n$
	from the closed disk centered at $(0,0)$ with radius $n$;
	we can see that $\bigcup_nL_n =\RR^2\setdiff\ol U$.
	Furthermore, in polar coordinates,
	\[\int_{L_n} f = \int_{[0,2\pi]\times[1+1/n,n]} \frac{1}{r}
	= 2\pi(\ln(n) - \ln(1 + 1/n)) = 2\pi\cdot\ln(n^2/(n+1)).\]
	Again this is unbounded as $n$ grows large,
	so $f$ is not integrable over $\RR^2\setdiff\ol U$.
\end{soln}

\begin{problem}
	Let $B$ be the portion of the first quadrant in $\RR^2$ lying
	between the hyperbolas $xy=1$ and $xy=2$ and the two lines
	$y=x$ and $y=4x$.
	Evalute $\int_Bx^2y^3$.
\end{problem}

\begin{soln}
	Transform $x\mapsto u/v$, $y\mapsto uv$ with $u,v>0$, so that
	the region bounded by the hyperbolas and lines maps to
	the rectangle bounded between $u = 1$, $u = \sqrt2$,
	$v = 1$, and $v = 2$. Then the Jacobian for the transformation is
	\[\begin{vmatrix}
		1/v&-u/v^2\\v&u
	\end{vmatrix} = 2u/v.\]
	Then $x^2y^3 = u^5/v$, so, by Fubini's theorem, the integral equals
	\[\int_R x^2y^3 = 2\int_1^2 \frac{1}{v^2}\cdot \int_1^{\sqrt2} u^6
	= 2\cdot\frac{1}{2}\cdot \frac{\sqrt{2}^7 - 1}{7} = \frac{8\sqrt2-1}{7}.\]
\end{soln}

\begin{problem}
	\begin{enumerate}[(a)]
		\ii If $f\colon\RR^2\to\RR$ is of class $C^1$,
		show that $f$ is not one-to-one.
		\ii If $f\colon\RR\to\RR^2$ is of class $C^1$,
		show that $f$ does not carry $\RR$ onto $\RR^2$.
		In fact, show that $f(\RR)$ contains no open set in $\RR^2$.
	\end{enumerate}
\end{problem}

\begin{soln}
	(a) If $\mathrm df = 0$ everywhere then $f$ is constant.
	Otherwise let $p = (p_1, p_2)$ be a point at which (WLOG)
	$\frac{\partial f}{\partial y}(p)\ne 0$.
	Then, by the implicit function theorem,
	there exists a neighborhood $U\subseteq\RR$ of $p_1$ and $V\subseteq\RR$
	of $p_2$ such that there is a $C^1$ function $g\colon U\to V$
	satisfying $g(p_1) = p_2$ such that $f(x, g(x)) = f(p)$
	for all $x\in U\times V$, proving that $f$ isn't injective.

	(b) We show that $f(\RR)$ has measure zero in $\RR^2$.
	
	Restrict $f$ first to an interval $[a,b]$.
	Then, if $f(x) = (f_1(x), f_2(x))$ for
	$C^1$ functions $f_1$ and $f_2$, by the mean value theorem,
	$f(x) - f(y) = (f_1'(\alpha_1), f_2'(\alpha_2))(x - y)$
	for some $\alpha_1,\alpha_2\in[a,b]$. Since $[a,b]$ is compact,
	$f_1'$ and $f_2'$, being continuous, are bounded on $[a,b]$,
	hence $\frac{\|f(x) - f(y)\|}{|x-y|}$ is bounded for $x,y\in[a,b]$.
	This implies $f$ is Lipschitz on $[a,b]$,
	so say $\|f(x)-f(y)\|\le L|x-y|$ for all $x,y\in[a,b]$.
	
	Let the points $a = t_0 \le t_1\le\cdots\le t_n=b$
	partition $[a,b]$ into intervals $[t_i, t_{i+1}]$ of equal length.
	Then, each set $f([t_i, t_{i+1}])$ can be covered with
	a cube of side length $L\cdot\frac{b-a}{2n}$, giving a covering
	of $f([a,b])$ with boxes of total volume at most
	\[n\cdot\left(L\cdot\frac{b-a}{2n}\right)^d =
	\left(\frac{L\cdot(b-a)}{2}\right)^d\cdot n^{1-d},\]
	which goes to zero as $n$ grows arbitrarily large,
	implying that $f([a,b])$ has measure zero.
	
	Now $f(\RR) = \bigcup_{n=1}^\infty f([-n,n])$,
	which is a union of countably many sets of measure zero,
	hence also has measure zero.
	
	This implies the result because all open sets do not have measure zero,
	on account of any open set containing an open ball,
	which clearly has nonzero measure.
\end{soln}
























\end{document}